\chapter{L'optimisation A2 : comptage, conception, hypothèses et limites}
\label{ch:a2}

Ce chapitre répond à QR1 et décrit la réalisation (QR2). Les propositions~\ref{prop:borne} à~\ref{prop:real} sont des énoncés mathématiques \emph{dans un modèle}; aucune conclusion de performance n'en est déduite (section~\ref{sec:perfborne}).

\section{Modèle de comptage}
\begin{definition}[Données]
Un \emph{$m$-vec} est un vecteur $v\in\mathbb{F}_{16}^m$, $m$ pair, stocké sur $b=m/2$ octets : $v_{2i}$ est le quartet bas de l'octet $i$, $v_{2i+1}$ le quartet haut. Les $K$ scalaires $c_0,\dots,c_{K-1}$ sont regroupés en $P=\lceil K/2\rceil$ paires (complétées par 0 si $K$ est impair). Pour chaque $k<K$ on calcule $\mathrm{acc}_k\leftarrow\mathrm{acc}_k\oplus c_k v$.
\end{definition}
Pour la paire $p$, $T_p[x]=(c_{2p}x)\oplus((c_{2p+1}x)\ll4)$ ; \vpshufb{} de largeur $W\in\{16,32\}$ octets effectue $W$ recherches indépendantes, chacune livrant, pour \emph{un} élément de $v$, les deux produits de la paire.

\begin{definition}[Classe $\mathcal{C}$]
Algorithmes qui (i) obtiennent chaque produit $c_kv_j$ par une recherche dans une table de 16 octets ; (ii) fournissent à chaque recherche un quartet isolé dans un octet (bit~7 nul) ; (iii) ne font pas deux recherches pour le même couple (élément, paire) hors du recouvrement entre registres ; (iv) utilisent librement XOR, ET, décalages, permutations, chargements. Sont \emph{hors} de $\mathcal{C}$ : GFNI, représentations \emph{bitsliced}, tables de 256 octets.
\end{definition}
La quantité comptée est $N(\text{algo},m)$ : nombre de \vpshufb{} exécutés \emph{par paire de scalaires et par $m$-vec}.

\section{Énoncés}
\begin{proposition}[Borne]\label{prop:borne}
Tout algorithme de $\mathcal{C}$ vérifie $N\ge\lceil m/32\rceil$.
\end{proposition}
\begin{proof}
Chaque élément $v_j$ ($j<m$) exige au moins une recherche par paire ; une instruction en fournit au plus $W\le32$. Donc $32N\ge m$.
\end{proof}

\begin{proposition}[Schéma officiel avec recouvrement]\label{prop:off}
Le schéma qui couvre les $b$ octets par $\lceil b/32\rceil$ chargements de 32 octets, le dernier reculé, et traite chaque registre par deux \vpshufb{} (impairs, pairs) a $N_{\mathrm{off}}(m)=2\lceil m/64\rceil$.
\end{proposition}
\begin{proof}
$\lceil b/32\rceil=\lceil m/64\rceil$ chargements, chacun fournissant 32 indices impairs et 32 pairs, soit deux \vpshufb{}.
\end{proof}

\begin{proposition}[Réalisation de la borne si $0<r\le32$]\label{prop:real}
Écrivons $m=64q+r$, $0\le r<64$, et $t=r/2$ le nombre d'octets de la queue. Si $0<r\le32$, l'algorithme suivant est dans $\mathcal{C}$ et a $N=2q+1=\lceil m/32\rceil$ :
\begin{enumerate}[nosep]
\item pour $g<q$ : charger l'intervalle $[32g,32g+32)$ et le traiter par deux \vpshufb{} ymm ;
\item pour la queue : charger les $t$ octets de $[32q,32q+t)$, construire $X=[\text{quartets bas des }t\text{ octets}\mid\text{quartets hauts des }t\text{ octets}]$ ($r\le32$ indices dans un seul registre) et appliquer \emph{un} \vpshufb{} ;
\item réassembler les sorties de la queue (même formule qu'un registre complet, restreinte aux $r$ positions).
\end{enumerate}
\end{proposition}
\begin{proof}
Le décompte est $2q+1$, et $\lceil m/32\rceil=2q+\lceil r/32\rceil=2q+1$ pour $0<r\le32$.
\end{proof}
\emph{Conditions de réalisation} (à satisfaire par le code, non par le modèle) : (a) le chargement de la queue ne lit que des octets de l'entrée ; (b) si $r\le16$ ($t\le8$), $X$ tient dans un registre 128 bits et la table utilisée est la \emph{voie basse} de $T_p$ ; si $16<r\le32$, $X$ occupe un registre 256 bits (voie basse = impairs, haute = pairs), ce qui exige l'égalité des deux voies ; (c) le surcoût de préparation de la queue est payé une fois par entrée, celui du réassemblage une fois par paire.

\begin{corollaire}[Limite exacte]
Pour $m\le192$, $\Delta(m):=N_{\mathrm{off}}(m)-\lceil m/32\rceil=1$ si $0<(m\bmod64)\le32$ et $0$ sinon. Dans le second cas, le code officiel atteint déjà la borne : aucun algorithme de $\mathcal{C}$ ne peut faire mieux en nombre de \vpshufb{}.
\end{corollaire}

\begin{table}[ht]\centering\small
\caption{Application aux paramètres actuels.}\label{tab:delta}
\begin{tabular}{@{}lrrrrrrl@{}}\toprule
 & $m$ & $q$ & $r$ & $N_{\mathrm{off}}$ & $\lceil m/32\rceil$ & $\Delta$ & Statut\\\midrule
MAYO-1 & 80 & 1 & 16 & 4 & 3 & 1 & implémentée (A2), queue de 8 octets\\
MAYO-2 & 64 & 1 & 0 & 2 & 2 & 0 & sans objet\\
MAYO-3 & 108 & 1 & 44 & 4 & 4 & 0 & impossible dans $\mathcal{C}$\\
MAYO-5 & 142 & 2 & 14 & 6 & 5 & 1 & implémentée (A2), queue de 7 octets + 1 de bourrage\\\bottomrule
\end{tabular}
\end{table}

\paragraph{Ce qui est implémenté et ce qui ne l'est pas.} Le prototype traite les queues d'au plus 16 éléments ($t\le8$, $r\le16$). La variante pour $16<r\le32$ \textbf{n'est ni implémentée ni testée} ; elle n'est pas un résultat acquis. Aucun des quatre paramètres actuels n'a $r\in(16,32]$.

\section{Conception d'A2}
\subsection{Structure du noyau}
Pour une entrée de $P_1$ et une ligne d'accumulateurs de $K$ vecteurs, A2 conserve inchangés les $q=\lfloor m/64\rfloor$ registres ymm complets sur les $32q$ premiers octets, puis traite la queue ainsi :
\begin{enumerate}[nosep]
\item \textbf{Préparation (une fois par entrée)} : chargement de 8 octets (\code{_mm_loadl_epi64}), décalage de 4 bits (\code{vpsrlw}), entrelacement \code{vpunpcklqdq} (quartets bas dans les octets 0--7, hauts dans 8--15), ET avec \code{0x0f} : quatre instructions (cinq dans la variante SAFE, qui ajoute un décalage d'octets).
\item \textbf{Produits} : un \code{vpshufb} xmm avec la \emph{moitié basse} de la table de multiplication déjà chargée, et un \code{vpxor} xmm dans un accumulateur xmm.
\item \textbf{Réassemblage (fin de ligne)} : séparation impair/pair identique à la partie ymm, et lecture-modification-écriture d'\emph{un} limbe de 64 bits (le limbe d'indice $4q=L-1$, soit 4 pour MAYO-1 et 8 pour MAYO-5).
\end{enumerate}
Dans l'assembleur de GCC~16.2.0, la boucle interne passe donc, par entrée, de 20 à 15 \code{vpshufb} pour MAYO-1 (10 ymm + 5 xmm) et de 36 à 30 pour MAYO-5 (24 ymm + 6 xmm).

\subsection{Variante par défaut et variante SAFE (MAYO-5)}
Pour MAYO-1, $b=40=32+8$ : la queue est exactement la fin de l'entrée, $[32,40)$. Pour MAYO-5, $b=71$ et $64+8=72$ :
\begin{itemize}[nosep]
\item \emph{défaut} : lit $[64,72)$, soit \textbf{un octet au-delà} de l'entrée (comme le noyau officiel, dont le dernier chargement couvre $[40,72)$) ;
\item \emph{SAFE} : lit $[63,71)$ puis décale le registre de 8 bits vers les poids faibles ; l'octet 63 est éliminé, les octets 64--70 sont placés aux positions 0--6, la position 7 reçoit 0. \textbf{Toute lecture reste dans l'entrée.}
\end{itemize}
La variante retenue pour une éventuelle intégration est SAFE seule. Les quartets de bourrage (positions 142 et 143 de MAYO-5, car $16L=144$) n'ont pas d'influence sur les positions utiles (section~\ref{sec:bourrage}).

\subsection{Préconditions}
(P1) $P_1$ est dans la disposition empaquetée, sans bourrage entre entrées, à l'adresse $P_1+e\cdot b$ ; (P2) les accumulateurs ont $L$ limbes avec bourrage de quartets seulement en fin de dernier limbe ; (P3) les tables de multiplication sont construites par la routine officielle (voies égales, lemme~\ref{lem:voies}) ; (P4) $m\in\{80,142\}$ : les noyaux A2 ne sont pas actifs pour les autres paramètres (condition de compilation \code{MAYO_TAIL_XMM_ACTIVE}) ; (P5) compilation sans \code{ENABLE_PARAMS_DYNAMIC}.

\section{La borne n'implique pas l'accélération}\label{sec:perfborne}
Une réduction du nombre de \vpshufb{} ne garantit pas une réduction du temps : s'y ajoutent le coût de préparation de la queue, les dépendances, les rechargements depuis la pile, les chargements non alignés et le débit des unités. Si seule la contrainte de débit du \vpshufb{} (un par cycle sur Skylake~\cite{nibbling}) bornait la boucle, le rapport de temps attendu serait $N_{A2}/N_{\mathrm{off}}$ : $15/20=0{,}750$ (MAYO-1) et $30/36=0{,}833$ (MAYO-5).

\begin{table}[ht]\centering\small
\caption{Rapport attendu par le seul débit de \vpshufb{} et rapport mesuré (S1, A2/officiel, noyaux isolés).}\label{tab:borne-mesure}
\begin{tabular}{@{}lccl@{}}\toprule
Noyau & Borne « débit seul » & Mesuré (S1) & Lecture\\\midrule
MAYO-1 \code{P1_times_Vt} & 0,750 & 0,780 & plus lent que la borne\\
MAYO-1 \code{P1t_times_Vt} & 0,750 & 0,710 & plus rapide que la borne\\
MAYO-5 \code{P1_times_Vt} & 0,833 & 0,807 & plus rapide que la borne\\
MAYO-5 \code{P1t_times_Vt} & 0,833 & 0,880 & plus lent que la borne\\\bottomrule
\end{tabular}
\end{table}
Deux cas sur quatre sont plus rapides que ne le permet le seul comptage de \vpshufb{} : d'autres mécanismes contribuent (par exemple moins de rechargements de pile : 11 $\to$ 5 pour MAYO-1 \code{P1_times_Vt}, 30 $\to$ 24 pour MAYO-5), sans que cela soit établi. Le modèle de débit n'explique donc pas à lui seul les gains. Une hypothèse de performance H-P (« sur une microarchitecture où le débit des shuffles est le goulot, le temps est borné inférieurement par $N\cdot P$ cycles par entrée ») ne peut se tester qu'avec des cycles cœur, pas des ticks TSC ; elle est postérieure à ce travail et reste ouverte.
